
\label{chap:Intro}


\cleardoublepage

\setcounter{section}{-1}
\section{Introducción}

En la última década, la transformación digital ha dejado de ser tendencia para convertirse en una necesidad estratégica para las organizaciones. Este proceso involucra la adopción de tecnologías emergentes para rediseñar operaciones, optimizar recursos y mejorar la adaptación frente a entornos cada vez más volátiles. En el contexto de la Cuarta Revolución Industrial (Industria 4.0), la convergencia de tecnologías como el Internet Industrial de las Cosas (IIoT), la inteligencia artificial (IA), la computación en la nube y la analítica avanzada está redefiniendo los modelos productivos y de gestión en todos los niveles de las organizaciones.



Según el Foro Económico Mundial, las empresas que adoptan estas tecnologías no solo logran mejoras en su eficiencia operativa, sino que también fortalecen su capacidad de respuesta ante crisis, interrupciones logísticas o cambios bruscos en la demanda global \cite{wef1}. En este sentido, la transformación digital se convierte en pieza clave para habilitar la resiliencia organizacional. Una evidencia particular es presentada en el artículo \cite{scd1}, donde se analizaron datos empresariales para evaluar el impacto de la transformación digital sobre la resiliencia de las cadenas de suministro. Los resultados muestran que la digitalización no solo tiene un efecto directo positivo en la capacidad de adaptación organizacional, sino que además fortalece la transparencia en la cadena de suministro, lo que contribuye a una mayor estabilidad frente a disrupciones externas. Aunque América Latina enfrenta desafíos en infraestructura digital y formación de talento, informes de la CEPAL subrayan que las empresas que han avanzado en procesos de transformación digital en esta región muestran una mayor capacidad de recuperación económica y adaptación estratégica frente a entornos cambiantes \cite{cepaL}.

En paralelo, el avance de la Industria 4.0 también plantea retos en el ámbito laboral. Según el Foro Económico Mundial, la implementación de tecnologías como la robótica, la automatización y la IA está transformando no solo los procesos productivos, sino también las funciones de los trabajadores dentro de las organizaciones. Esta transformación exige una actualización constante de habilidades técnicas y digitales, así como una reconfiguración de los roles laborales tradicionales para adaptarse a entornos cada vez más conectados \cite{wef2}.

A pesar de estos avances, muchas industrias aún operan con arquitecturas tradicionales, como la ISA-95, que, aunque están ampliamente consolidadas, pueden limitar la interoperabilidad y dificultar la integración con nuevas tecnologías. Frente a esto, la arquitectura Unified Namespace (UNS) emerge como una alternativa para centralizar, estandarizar y estructurar el flujo de datos en tiempo real. Esta arquitectura facilita la integración entre sistemas, elimina silos de información y permite el despliegue eficiente de soluciones avanzadas en analítica de datos.


\subsection{\centering Justificación}

La transformación digital es un factor determinante en la evolución de los entornos industriales hacia modelos más inteligentes, flexibles y resilientes. La ndustria 4.0 ha introducido tecnologías que permiten una automatización avanzada y una integración más profunda entre los sistemas físicos y digitales. No obstante, aún persisten desafíos importantes en cuanto a la forma en que las empresas acceden, estructuran y explotan eficientemente la información generada en sus procesos operativos.


Uno de los principales obstáculos radica en la fragmentación y dispersión de los datos, tanto estructurados como no estructurados. El UNS surge como una solución para centralizar y estandarizar el flujo de datos en tiempo real. Aunque implementar una arquitectura UNS mejora la accesibilidad y la integración de los datos, por sí sola no garantiza el aprovechamiento inteligente de la información. Para lograr un verdadero valor agregado, es necesario incorporar herramientas que permitan interpretar los datos de forma eficiente.


En este sentido, los Modelos de Lenguaje de Gran Tamaño (LLMs) ofrecen un potencial transformador. Gracias a su capacidad para comprender lenguaje natural, generar texto coherente y realizar búsquedas semánticas sobre grandes volúmenes de información, los LLMs pueden actuar como una capa de inteligencia contextual sobre la arquitectura UNS. Esta sinergia habilita nuevas formas de interacción entre los operadores y los sistemas industriales, permitiendo, por ejemplo, responder consultas técnicas mediante lenguaje natural, generar reportes automatizados y acceder a conocimiento distribuido sin necesidad de navegar por múltiples fuentes.


\subsection{\centering Descripción del proyecto}
En el contexto industrial actual, las empresas generan y consumen grandes volúmenes de información provenientes de fuentes heterogéneas. Esta información abarca tanto datos estructurados (registros de sistemas ERP, MES, SCADA, y bases de datos) como datos no estructurados (manuales técnicos, reportes operativos, y documentos de mantenimiento). Uno de los principales desafíos es la fragmentación, dispersión y falta de estandarización de estos datos. En este escenario, la arquitectura UNS surge como una solución para centralizar, estandarizar y visualizar de manera unificada la información de todos los niveles de la operación en tiempo real. No obstante, la simple centralización no garantiza el aprovechamiento inteligente de la información. Se requieren herramientas que permitan realizar análisis semánticos, búsquedas contextuales y generación automatizada de conocimiento útil. Los LLMs ofrecen un potencial transformador gracias a su capacidad para comprender lenguaje natural, extraer significado de textos no estructurados y generar respuestas contextualizadas, por lo que tienen la capacidad de actuar como una capa de ``inteligencia'' sobre la arquitectura UNS. El presente proyecto de investigación tiene como propósito proponer un mecanismo para la integración efectiva de herramientas basadas en Modelos de LLMs dentro de arquitecturas UNS aplicadas al contexto de la Industria 4.0. La investigación busca determinar cómo la incorporación de estas herramientas puede potenciar las capacidades del UNS, permitiendo funcionalidades avanzadas como búsquedas inteligentes, consultas semánticas, generación automática y dinámica de reportes en entornos industriales reales.


\subsection{\centering Objetivos}
A continuación, se presentan el objetivo general y los objetivos específicos del trabajo de investigación.


\textbf{Objetivo general}

Integrar una herramienta basada en modelos de lenguaje de gran tamaño para la analítica de datos de una arquitectura Unified Namespace aplicada a un caso de industria 4.0.

\vspace{0.5cm}


\textbf{\centering Objetivos específicos}
\begin{itemize}
    \item \textbf{OE1:} Caracterizar un caso de estudio en Industria 4.0, identificando los requerimientos y desafíos específicos de analítica de datos.

    \item \textbf{OE2:} Diseñar e implementar una arquitectura Unified Namespace, estableciendo la estructura de datos.

    \item \textbf{OE3:} Seleccionar una herramienta basada en modelo de lenguaje de gran tamaño, considerando criterios de adaptabilidad al caso de estudio.

    \item \textbf{OE4:} Integrar la herramienta de lenguaje de gran tamaño a la arquitectura Unified Namespace, proporcionando al modelo el contexto correspondiente al caso de estudio.

    \item \textbf{OE5:} Validar el funcionamiento del sistema integrado mediante métricas de desempeño.

\end{itemize}

\subsection{\centering Esquema del trabajo de investigación}

Este trabajo de investigación está organizado de la siguiente manera:

\color{blue}
\begin{itemize}
    \item  \textbf{Capítulo~\ref{cap:caso_estudio}:} Caracterización del caso de estudio de Industria 4.0
    \item  \textbf{Capítulo~\ref{tecnicasdiseñouns}:} Implementación de arquitectura UNS
    \item  \textbf{Capítulo~\ref{seleccionllm}:} Selección del mecanismo de integración y herramienta LLM
    \item  \textbf{Capítulo~\ref{implementacionintegracion}:} Integración de la herramienta LLM a la arquitectura UNS
    \item  \textbf{Capítulo~\ref{validacionsistema}:} Validación del sistema integrado

\end{itemize}
\color{black}